% Generated draft: requires mathematical and source review.
\documentclass[11pt]{article}
\usepackage[margin=25mm]{geometry}
\usepackage{fontspec}
\setmainfont{Latin Modern Roman}
\usepackage{amsmath,amssymb,mathtools,unicode-math}
\setmathfont{Latin Modern Math}
\usepackage{enumitem,needspace,xurl,hyperref}
\setlength{\parindent}{0pt}
\setlength{\parskip}{5pt}
\setlength{\emergencystretch}{3em}
\newcommand{\N}{\mathbb{N}}
\newcommand{\Z}{\mathbb{Z}}
\newcommand{\Q}{\mathbb{Q}}
\newcommand{\R}{\mathbb{R}}
\newcommand{\C}{\mathbb{C}}
\newcommand{\F}{\mathbb{F}}
\newcommand{\A}{\mathbb{A}}
\newcommand{\PP}{\mathbb{P}}
\newcommand{\E}{\mathbb{E}}
\newcommand{\eps}{\varepsilon}
\newcommand{\dd}{\,\mathrm{d}}
\begin{document}
\section*{UnsolvedMath ID: 30003141}
\textbf{Source problem code:} catalogue entry
\subsection{Definitions and mathematical statement}
Let \(g\in\Z\) with \(g\geq5\). In the conventional interpretation, \(M_g\) denotes the moduli space or moduli stack of smooth, proper, connected curves of genus \(g\). The supplied record does not specify the base field or whether the stack or coarse moduli space is intended. It also does not define the coefficient object \(V_{2222}\), the cohomology theory and coefficient field, or the meaning of the notation \(\operatorname{RH}^i(M_g,V_{2222})\). These symbols are therefore retained as source notation, without imposing additional conventions.

The source asks whether Faber's conjecture fails for genera \(g\geq5\). It does not specify the version of Faber's conjecture, the pairing whose failure is intended, or whether the question asks for a counterexample in at least one such genus or for failure in every such genus. Consequently, the general question cannot be made into an unambiguous quantified conjecture from the supplied record alone.

The concrete genus-five target recorded by the source is to establish both assertions
\[
\operatorname{RH}^{4}(M_5,V_{2222})\neq0,
\qquad
\operatorname{RH}^{2}(M_5,V_{2222})=0.
\]
The second assertion is described as expected, rather than established. According to the supplied record, proving these assertions together would produce a counterexample to the intended form of Faber's conjecture. The record does not supply the definitions or the precise pairing needed to formulate or justify that implication independently. This is therefore a statement-only transcription with unresolved mathematical specifications, not an assertion that either cohomology calculation or the claimed counterexample has been proved.
\subsection{Short English statement}
Investigate the proposed failure of Faber's conjecture in genus at least five. The supplied genus-five test asks for nonvanishing of \(\operatorname{RH}^{4}(M_5,V_{2222})\) together with the expected vanishing of \(\operatorname{RH}^{2}(M_5,V_{2222})\). The coefficient object, cohomology notation, precise conjecture, pairing, and scope across genera require clarification before this becomes a fully specified mathematical problem.
\subsection{Sources}
\begin{itemize}
\item Supplied local UnsolvedMath registry record, permanent ID \(30003141\), source problem code \texttt{catalogue entry}, entitled \emph{Failure of Faber's Conjecture in Genus Five}, in the Oberwolfach Reports Open Problems collection. No individual proposer is identified; the supplied proposed year is 2016.
\item Supplied source citation: \emph{Moduli spaces and Modular forms (2017)}, Oberwolfach Reports, pp.\ 1259--1318, DOI \texttt{10.4171/owr/2016/23}. Original OWR report: \url{https://doi.org/10.4171/owr/2016/23}. The record describes Oberwolfach Reports as documenting mathematical workshops and their open-problem discussions.
\item Supplied literature source: Pixton, \emph{The tautological ring}, \url{https://math.mit.edu/~apixton/papers/tautological.pdf}.
\item Supplied literature assessment dated exactly \texttt{2026-08-21}: the record labels the problem open and reports only numerical evidence for a possible genus-five counterexample; its assessment states that no proof of the required nonzero cohomology group or failure of the stated pairing was located. This is supplied dated curation derived from the OWR record, not an independently verified literature assessment. The record requests independent verification of time-sensitive status and cited claims.
\end{itemize}
\end{document}
